\subsection{A coherent unit cofactor in the fixed pure-three norm tower}

\begin{statement}{A.138}{Proved}
{An actual unit norm construction and its pro-three compatible half-powers}
Choose positive roots
\[
 \alpha_n=3^{1/3^n},\qquad
 K_n=\Q(\alpha_n),\qquad
 L_n=\Q(\alpha_n,\mu_{3^n}),\qquad
 \alpha_{n+1}^3=\alpha_n .
\tag{A.138.1}
\]
Every norm
\(\mathcal O_{L_{n+1}}^\times\to
\mathcal O_{L_n}^\times\) is surjective.
There exist actual global units \(v_n\in\mathcal O_{K_n}^\times\)
with
\[
 \begin{gathered}
 N_{K_{n+1}/K_n}(v_{n+1})=v_n,\qquad
 N_{K_n/\Q}(v_n)=1,\\
 v_1=(\alpha_1^2-2)^2 .
 \end{gathered}
\tag{A.138.2}
\]
Write
\(\widehat{\mathcal O_{K_n}^\times}_{(3)}
=\varprojlim_r\mathcal O_{K_n}^\times/
(\mathcal O_{K_n}^\times)^{3^r}\).
The unique half-powers \(U_n=v_n^{1/2}\) in these
pro-three completions give a single compatible family
\[
 \begin{gathered}
 U_1=\alpha_1^2-2,\qquad
 N_{K_{n+1}/K_n}(U_{n+1})=U_n,\\
 \beta_n=(1-\alpha_n)U_n
       \in\widehat{K_n^\times}_{(3)},\qquad
 N_{K_{n+1}/K_n}(\beta_{n+1})=\beta_n,\qquad
 N_{K_n/\Q}(\beta_n)=-2 .
 \end{gathered}
\tag{A.138.3}
\]
The completed unit factors have zero finite valuations,
so \(\beta_n\) has precisely the divisor of \(1-\alpha_n\),
including its canonical divisor above two.

For the actual norm-one torus
\(T_1=R^1_{K_1/\Q}\mathbf G_m\), the quotient
\[
 \lambda:T_1[3]\longrightarrow\F_3,\qquad
        (z_0,z_1,z_2)\longmapsto z_1+2z_2
\tag{A.138.4}
\]
in the labeled Kummer coordinates satisfies
\[
 \lambda_*\kappa_{T_1}(\alpha_1^2-2)=\chi_9 .
\tag{A.138.5}
\]
Here \(\chi_9:G_\Q\to\F_3\) is the cyclic cubic
cyclotomic character normalized by \(\chi_9(4)=1\)
on \((\Z/9)^\times/\{\pm1\}\).
Thus the nonzero first-depth quotient in (A.138.5)
occurs in a coherent family on the same tower.
The tower in (A.138.1) remains fixed; this construction
does not supply a norm family after changing its
radicand to an infinite-support completed coordinate.
\end{statement}
\begin{proof}
\emph{The actual field degrees and the single ramified prime.}
The polynomial \(X^3-3\) is Eisenstein over \(\Q_3\).
Its root field has degree three and is not Galois:
a normal root field would contain \(\mu_3\), whose
degree over \(\Q_3\) is two.
No abelian extension of \(\Q_3\) can contain this
non-Galois subfield.
In particular 3 is not a cube in any local cyclotomic
field \(\Q_3(\mu_{3^j})\).
Kummer theory therefore gives degree \(3^n\) for
adjoining \(\alpha_n\) to that field whenever \(j\ge n\).
The actual automorphisms
\[
 \tau(\alpha_n)=\zeta_{3^n}\alpha_n,\qquad
 \sigma_u(\alpha_n)=\alpha_n,\qquad
 \sigma_u(\zeta_{3^n})=\zeta_{3^n}^{\,u}
\tag{A.138.6}
\]
give the local and global holomorph, both of order
\(3^n\varphi(3^n)\).
The degree assertion also follows from the valuation
and pure-root argument in A.105, with the radicand
fixed here to be exactly 3.
For \(u=-1\), the commutator action on the
translation subgroup is multiplication by \(-2\),
which is invertible modulo \(3^n\).
Consequently the maximal abelian subextension of
the local field \(L_{n,3}/\Q_3\) is its cyclotomic
subfield.
An unramified quotient of a local Galois group is
cyclic and hence factors through this subfield.
The cyclotomic subfield is totally ramified, so
that quotient is trivial.
Thus \(L_n/\Q\) has exactly one prime over three
and is totally ramified there.
The polynomial and cyclotomic discriminants show
that no other finite prime ramifies.

Put \(k=\Q(\mu_3)\).
The same statements apply to
\[
 F_n=L_n(\mu_{3^{n+1}}),\qquad
 [F_n:L_n]=3,\qquad [L_{n+1}:F_n]=3 .
\tag{A.138.7}
\]
Both steps are cyclic, the first by its cyclotomic
action and the second by its Kummer action.
The fields \(L_n\) and \(F_n\) are finite Galois
three-extensions of \(k\), unramified outside
its unique prime above three.

\emph{The precise three-class-group argument.}
The discriminant of \(k\) is \(-3\).
Minkowski's bound
\cite[Theorem 4.3]{AgentMilneANT}
places an integral representative of every ideal
class below
\((2/\pi)\sqrt3<2\).
Its ideal norm must therefore be one, so
\(\operatorname{Cl}(k)=0\).

In fact every finite Galois three-extension
\(D/k\) unramified outside that one prime is
totally ramified there.
Let \(G=\operatorname{Gal}(D/k)\) and let
\(I\subset G\) be the inertia group at one prime
above it.
The quotient by the normal closure of \(I\)
is unramified at every finite prime.
If that quotient were nontrivial, its finite
three-group would have a nontrivial cyclic
quotient, giving a nontrivial unramified
abelian three-extension of \(k\), contrary
to its class number one.
Thus \(I\) normally generates \(G\).
A proper subgroup of a finite three-group
lies in a maximal subgroup, and every maximal
subgroup is normal of index three.
Such a proper subgroup cannot normally generate.
Hence \(I=G\), proving total ramification.

Apply this to the Hilbert three-class field
\(D\) of \(L_n\).
It is normal over \(k\), because every
\(k\)-automorphism of \(L_n\) preserves its
maximal unramified abelian three-extension.
Its Galois group over \(k\) is a three-group,
and no finite prime except three ramifies.
The preceding argument makes \(D/k\) totally
ramified at three.
But \(D/L_n\) is unramified, so inertia in
\(\operatorname{Gal}(D/k)\) has trivial
intersection with \(\operatorname{Gal}(D/L_n)\).
Since inertia is the entire group,
that intersection can be trivial only when
\(D=L_n\).
The identical argument applies to \(F_n\).
We have proved
\[
             \operatorname{Cl}(L_n)\{3\}=0,
             \qquad \operatorname{Cl}(F_n)\{3\}=0 .
\tag{A.138.8}
\]

\emph{Every cyclic cubic step has a surjective actual unit norm.}
Let \(E/F\) be either step in (A.138.7), and
let \(\varepsilon\in\mathcal O_F^\times\).
At an unramified finite place it is a local
norm: norm maps on units of unramified local
extensions are onto, including split places.
Both fields are totally imaginary.
The cyclic algebra associated to \(E/F\) and
\(\varepsilon\) therefore has zero local invariant
everywhere except possibly at the single prime
over three.
Global reciprocity makes that last invariant
zero as well.
The Hasse norm theorem for the cyclic extension
\cite[V, Theorem 5.3 and VIII, Theorem 3.1]{MilneCFT}
gives \(\beta\in E^\times\) with
\(N_{E/F}\beta=\varepsilon\).

Let \(\sigma\) generate \(\operatorname{Gal}(E/F)\).
The divisor \(D=\operatorname{div}(\beta)\)
has ideal norm zero.
On each orbit of primes of \(E\), ideal norm
is the common residue degree times the sum
of the coefficients.
For an orbit of size three, the kernel of that
sum is exactly the image of \(\sigma-1\) on
the free integral prime module.
For an orbit of size one, both kernels are zero.
It follows, with actual finite-support divisors,
that \(D=(\sigma-1)J\) for a fractional ideal
divisor \(J\).
By (A.138.8) the order \(d\) of \([J]\) is
prime to three.
Choose \(\gamma\in E^\times\) with
\(\operatorname{div}(\gamma)=dJ\).
The element
\[
 \varepsilon'=\beta^d\frac{\gamma}{\sigma\gamma}
       \in\mathcal O_E^\times,\qquad
 N_{E/F}(\varepsilon')=\varepsilon^d
\tag{A.138.9}
\]
is an actual unit.
For integers \(a,b\) with \(ad+3b=1\), the
unit \((\varepsilon')^a\varepsilon^b\) has
norm \(\varepsilon\).
This proves actual unit-norm surjectivity
at each step, and hence for \(L_{n+1}/L_n\).

\emph{One coherent family and its exact completion maps.}
Write \(U=\alpha_1^2-2\).
The identity
\[
 (\alpha_1^2-2)(2\alpha_1^2+3\alpha_1+4)=1,
                  \qquad N_{K_1/\Q}(U)=9-8=1
\tag{A.138.10}
\]
shows that it is a global unit of norm one.
Starting with the actual unit \(\eta_1=U\)
of \(L_1\), choose actual units
\(\eta_{n+1}\) with
\(N_{L_{n+1}/L_n}(\eta_{n+1})=\eta_n\).
The proved surjectivity permits every step.
Define
\[
                  v_n=N_{L_n/K_n}(\eta_n).
\tag{A.138.11}
\]
Both paths from \(L_{n+1}\) to \(K_n\)
give the same field norm.
Consequently
\[
 \begin{aligned}
 N_{K_{n+1}/K_n}(v_{n+1})
 &=N_{L_{n+1}/K_n}(\eta_{n+1})\\
 &=N_{L_n/K_n}(\eta_n)=v_n .
 \end{aligned}
\tag{A.138.12}
\]
At the first level \(L_1/K_1\) has degree two
and \(U\in K_1\), so \(v_1=U^2\).
Norm transitivity and (A.138.10) give
\(N_{K_n/\Q}(v_n)=1\).

For every \(r\), multiplication by two is
invertible on
\(\mathcal O_{K_n}^\times/
(\mathcal O_{K_n}^\times)^{3^r}\).
Its inverse commutes with the coefficient
transition maps, so it induces the unique
half-power map on the indicated pro-three
completion.
It also commutes with every field norm.
This proves all the assertions for \(U_n\).
The natural map
\[
 \widehat{\mathcal O_{K_n}^\times}_{(3)}
               \longrightarrow\widehat{K_n^\times}_{(3)}
\tag{A.138.13}
\]
is injective: at each finite exponent a root
of a global unit is a global unit, as its
finite valuations must vanish.
Thus these are specified elements of the
same completed multiplicative groups used
by the norm tower.
Finally,
\(N_{K_{n+1}/K_n}(1-\alpha_{n+1})=1-\alpha_n\)
and \(N_{K_n/\Q}(1-\alpha_n)=-2\).
They prove (A.138.3) and its divisor assertion.
No half-power in the full all-exponent
completion is used.

\emph{The actual first-depth quotient.}
Label the conjugates by
\(\alpha_i=\zeta_3^i\alpha_1\), \(i\in\Z/3\).
Then
\[
 T_1[3]=
 \{(z_0,z_1,z_2)\in\F_3^3:\ z_0+z_1+z_2=0\}
                  \otimes\chi_3 .
\tag{A.138.14}
\]
The affine action on the labels gives the
literal equivariance of \(\lambda\):
translation changes its formula by a multiple
of the zero coordinate sum, and multiplication
by \(u\in\{\pm1\}\) contributes \(u\) from
the labels and \(u\) from the cyclotomic
twist, whose product is one.
Thus its target is the trivial \(\F_3\)-module.
The Kummer class of the global unit \(U\)
is unramified away from three: at every
other prime its three root equations over
the unramified splitting algebra are finite
etale on units.
The same is true of its image under
\(\lambda\).

The only one-dimensional space of cubic
characters of \(\Q\) unramified outside
three is that of
\(\Q(\mu_9)^+/\Q\).
Indeed the ray class description of the
abelian extensions of \(\Q\)
\cite[I, Theorem 4.16 and V, Example 3.10]{MilneCFT}
puts them in the three-power cyclotomic tower;
the exponent-three quotient of
\(\Z_3^\times\) is
\((\Z/9)^\times/\{\pm1\}\).
It remains to determine the scalar.

At \(q=61\), take the prime of \(L_1\) with
\(\alpha_1=4\) and \(\zeta_3=47\) in the
residue field.
The three labels are \((4,5,52)\).
The number 2 is a primitive residue unit,
and \(2^{20}=47\).
The residues of the three units
\(\alpha_i^2-2\) are
\((14,23,18)\), whose logarithms to base 2
modulo three are \((2,0,1)\).
Thus their image under \(\lambda\) is 2.
On the other hand \(61\equiv7=4^2\pmod9\),
so \(\chi_9(\operatorname{Frob}_{61})=2\).
This proves (A.138.5) with the stated
normalization.

The norm diagram (A.138.12) keeps all the
fields in (A.138.1) fixed.
It proves a coherent extension of this
particular first-depth unit cofactor.
It contains no transition to a tower with
a changed completed radicand, which is
the remaining step in the ghost construction.
\end{proof}

The exact certificate
\path{code/constant_tower_unit_certificate.py}
checks the unit identity and norm, the
cyclic-prime divisor correction, the compatible
half-power coefficient maps, and the labeled
Frobenius computation at 61.
It does not claim to have enumerated the
infinite family of units chosen in the proof.
